\documentclass[11pt]{article}

\usepackage[margin=1in]{geometry}
\usepackage[T1]{fontenc}
\usepackage{lmodern}
\usepackage{microtype}
\usepackage{amsmath,amssymb,amsthm,mathtools}
\usepackage{mathrsfs}
\usepackage[dvipsnames]{xcolor}
\usepackage[most]{tcolorbox}
\usepackage{tikz}
\usetikzlibrary{decorations.markings,patterns}
\usepackage{enumitem}
\usepackage{tocloft}
\usepackage[hidelinks]{hyperref}

\definecolor{courseblue}{HTML}{1F4E79}
\definecolor{lightblue}{HTML}{EAF2F8}
\definecolor{lightgreen}{HTML}{EAF7EA}
\definecolor{lightpink}{HTML}{FCE8EF}
\definecolor{lightyellow}{HTML}{FFF8D6}
\definecolor{warningyellow}{HTML}{D6A000}
\definecolor{exercisebg}{HTML}{FDEAEA}
\definecolor{exerciseframe}{HTML}{B54040}
\definecolor{examplebg}{HTML}{F4FAFE}
\definecolor{exampleframe}{HTML}{85C1E9}
\definecolor{darkgray}{HTML}{333333}
\definecolor{lightorange}{HTML}{FDEBD8}
\definecolor{orangeframe}{HTML}{E08A3C}
\definecolor{remarkbg}{HTML}{F3EEFB}
\definecolor{remarkframe}{HTML}{9B7FC9}

\setlength{\parindent}{0pt}
\setlength{\parskip}{0.45em}
\setlist[itemize]{topsep=0.25em,itemsep=0.15em}
\renewcommand{\cftsecleader}{\cftdotfill{\cftdotsep}}
\renewcommand{\cftsubsecleader}{\cftdotfill{\cftdotsep}}

\theoremstyle{plain}
\newtheorem{theorem}{Theorem}[section]
\newtheorem{proposition}[theorem]{Proposition}
\newtheorem{lemma}[theorem]{Lemma}
\newtheorem{corollary}[theorem]{Corollary}
\theoremstyle{definition}
\newtheorem{definition}[theorem]{Definition}
\newtheorem{example}[theorem]{Example}
\newtheorem{exercise}[theorem]{Exercise}
\theoremstyle{remark}
\newtheorem*{remark}{Remark}
\newtheorem*{fact}{Fact}
\newtheorem*{warning}{Warning}

\tcolorboxenvironment{definition}{
  enhanced,
  unbreakable,
  colback=lightgreen,
  colframe=ForestGreen,
  boxrule=0.7pt,
  arc=2pt,
  left=6pt,
  right=6pt,
  top=5pt,
  bottom=5pt
}

\tcolorboxenvironment{remark}{
  enhanced,
  colback=remarkbg,
  colframe=remarkframe,
  before upper={\setlength{\parskip}{0.5em}},
  boxrule=0.7pt,
  arc=2pt,
  left=6pt,
  right=6pt,
  top=5pt,
  bottom=5pt
}

\tcolorboxenvironment{warning}{
  enhanced,
  colback=lightyellow,
  colframe=warningyellow,
  boxrule=0.7pt,
  arc=2pt,
  left=6pt,
  right=6pt,
  top=5pt,
  bottom=5pt
}

\tcolorboxenvironment{exercise}{
  enhanced,
  colback=exercisebg,
  colframe=exerciseframe,
  boxrule=1pt,
  arc=2pt,
  left=6pt,
  right=6pt,
  top=5pt,
  bottom=5pt
}

\tcolorboxenvironment{example}{
  enhanced,
  colback=examplebg,
  colframe=exampleframe,
  boxrule=0.7pt,
  arc=2pt,
  left=6pt,
  right=6pt,
  top=5pt,
  bottom=5pt
}

\numberwithin{equation}{section}

\newcommand{\be}{\begin{eqnarray}}
\newcommand{\ee}{\end{eqnarray}}
\newcommand{\C}{\mathbb{C}}
\newcommand{\R}{\mathbb{R}}
\newcommand{\N}{\mathbb{N}}
\newcommand{\Z}{\mathbb{Z}}
\DeclareMathOperator{\Log}{Log}
\DeclareMathOperator{\Arg}{Arg}

\DeclareMathOperator{\Hol}{\mathcal{H}\mathrm{ol}}
\DeclareMathOperator{\Real}{Re}
\DeclareMathOperator{\Imag}{Im}
\newcommand{\A}{\mathcal{A}}
\newcommand{\CT}{\C[[T]]}
\newcommand{\CLT}{\C((T))}
\tcbset{breakable}
% Compact spacing around diagrams; keep the main text at 11pt.
\newenvironment{lecturefigure}
 {\par\addvspace{0.25em}\begingroup\setlength{\parskip}{0pt}\centering}
 {\par\endgroup\addvspace{0.25em}}
\begin{document}

\begin{center}
  \colorbox{lightblue}{%
    \begin{minipage}{0.94\textwidth}
      \vspace{0.55em}
      \begin{tabular*}{\textwidth}{@{\extracolsep{\fill}}lr}
        {\color{courseblue}\bfseries\LARGE Lecture 22} &
        {\color{courseblue}\bfseries\LARGE MATH 411}\\[0.35em]
        \textbf{Fall 2026} & \textbf{Date: Oct 30, 2026}\\
        & \textbf{Instructor: Manish Patnaik}
      \end{tabular*}
      \vspace{0.45em}
    \end{minipage}}
\end{center}

\setcounter{tocdepth}{1}
\tableofcontents
\vspace{0.45em}

% Source: handwritten Lecture 21, Oct. 31, page 4, on isolated
% singularities. The Laurent expansion theorem, proof, and computational
% examples are in Lecture 21. The function-side characterizations below
% are proved in the next lecture.
In Lecture~21 we proved the Laurent expansion theorem and computed
expansions on different annuli and about different centres. Today we use
Laurent series to classify isolated singularities as removable
singularities, poles, or essential singularities.

\section{Isolated singularities}
% Handwritten Lecture 21, page 4. The function-side characterizations in
% the table are stated here; their proofs come in the next lecture.
\begin{definition}
Let $U$ be open and $z_0\in U$. A function $f$ holomorphic on
$U\setminus\{z_0\}$ is said to have an \emph{isolated singularity} at $z_0$.
\end{definition}
We are interested in functions $f$ on an open set that are holomorphic
except for, at worst, isolated singularities. If $f$ has an isolated
singularity at $z_0$, choose $\rho>0$ with $D(z_0,\rho)\subseteq U$. Then $f$
is holomorphic on the punctured disc $\{0<|z-z_0|<\rho\}$, an annulus with
inner radius $0$, so the open-annulus corollary of Lecture~21 gives a Laurent
expansion
\[
 f(z)=\sum_{n=-\infty}^{\infty}a_n(z-z_0)^n\qquad(0<|z-z_0|<\rho).
\]
By uniqueness of the Laurent expansion (Lecture~21), the coefficients depend only on $f$ near
$z_0$, not on $\rho$.

\begin{definition}
In this situation the singularity at $z_0$ is
\begin{enumerate}
 \item \emph{removable} if $a_n=0$ for all $n<0$;
 \item a \emph{pole} if $a_n\ne0$ for some, but only finitely many, $n<0$.
 If $a_{-m}\ne0$ and $a_n=0$ for $n<-m$, the pole has \emph{order} $m$;
 \item an \emph{essential singularity} if $a_n\ne0$ for infinitely many
 $n<0$.
\end{enumerate}
\end{definition}

The three types are also visible from the behaviour of $f$ near $z_0$:

\begin{center}
\renewcommand{\arraystretch}{1.5}
\begin{tabular}{|l|p{5.2cm}|p{5.4cm}|}
 \hline
 \textbf{Type} & \textbf{Behaviour of $f$ near $z_0$} &
 \textbf{Laurent expansion at $z_0$}\\
 \hline
 Removable & $f$ is bounded near $z_0$ &
 no negative terms:\newline $f=\sum_{n\geq0}a_n(z-z_0)^n$\\
 \hline
 Pole & $|f(z)|\to\infty$ as $z\to z_0$ &
 finitely many negative terms:\newline $f=\sum_{n\geq-m}a_n(z-z_0)^n$\\
 \hline
 Essential & the values of $f$ near $z_0$ are dense in $\C$ &
 infinitely many negative terms\\
 \hline
\end{tabular}
\end{center}
That the two columns match is a theorem, proved next lecture (it includes
Riemann's removable singularity theorem and the Casorati--Weierstrass
theorem).

\begin{example}
Some examples of each type:
\begin{itemize}
 \item $\dfrac{\sin z}{z}=1-\dfrac{z^2}{3!}+\dfrac{z^4}{5!}-\cdots$ has a
 removable singularity at $0$; setting its value at $0$ to be $1$ gives an
 entire function.
 \item $\dfrac{1}{z(z-1)}$ has poles of order $1$ at $0$ and at $1$, by
 the expansions computed in Lecture~21.
 \item $\dfrac{1}{z^2}$ has a pole of order $2$ at $0$.
 \item $e^{1/z}=\displaystyle\sum_{n\geq0}\frac{1}{n!}z^{-n}$ has an
 essential singularity at $0$.
\end{itemize}
\end{example}

\begin{exercise}
Prove the easy half of the table for the first two rows. (a) If the
Laurent expansion at $z_0$ has no negative terms, show that $f$ is bounded
near $z_0$ and extends holomorphically to $z_0$. (b) If $f$ has a pole of
order $m$ at $z_0$, write $f(z)=(z-z_0)^{-m}g(z)$ with $g$ holomorphic near
$z_0$ and $g(z_0)\ne0$, and deduce that $|f(z)|\to\infty$ as $z\to z_0$.
\end{exercise}

\end{document}
